\documentclass[11pt,a4paper]{article}
\usepackage[italian]{babel}
\usepackage[margin=2.5cm]{geometry}
\usepackage{parskip}
\usepackage{amsmath, amssymb, amsthm}
\usepackage{booktabs}
\usepackage{array}
\usepackage{xcolor}
\usepackage{needspace}
\usepackage{enumitem}
\usepackage{listings}
\usepackage{tikz}
\usepackage[most]{tcolorbox}
\usepackage[hidelinks]{hyperref}
\usetikzlibrary{decorations.pathreplacing, shapes.geometric, arrows.meta, positioning}

% Niente vedove/orfane: i paragrafi non lasciano righe isolate a fine/inizio pagina
\widowpenalty=10000
\clubpenalty=10000

% Elenchi più compatti
\setlist{itemsep=2pt, parsep=0pt, topsep=4pt, partopsep=0pt}

% --- Riquadri con bordo nero, senza numero (si spezzano tra due pagine) ---
% Uso: \begin{definizione} ... \end{definizione}
%      \begin{teorema}[Nome] ... \end{teorema}  (il nome è facoltativo)
\tcbset{nota/.style={enhanced jigsaw, breakable, colback=white,
  colframe=black, boxrule=0.5pt, sharp corners}}
\NewTColorBox{definizione}{}{nota,
  before upper={\textbf{Definizione.}\ }}
\NewTColorBox{teorema}{o}{nota,
  before upper={\textbf{Teorema\IfValueT{#1}{ (#1)}.}\ }}
\NewTColorBox{proposizione}{o}{nota,
  before upper={\textbf{Proposizione\IfValueT{#1}{ (#1)}.}\ }}
\NewTColorBox{proprieta}{o}{nota,
  before upper={\textbf{Proprietà\IfValueT{#1}{ (#1)}.}\ }}
\NewTColorBox{assioma}{o}{nota, boxrule=1pt,
  before upper={\textbf{Assioma\IfValueT{#1}{ (#1)}.}\ }}

% --- Ambienti semplici (senza riquadro) ---
% Il titolo sta in cima al blocco, anche se il contenuto inizia con un riquadro o
% con codice; e non resta mai da solo in fondo a una pagina.
% Uso: \begin{esempio}[Titolo facoltativo] ... \end{esempio}
\newcommand{\newrunin}[2]{%
  \NewDocumentEnvironment{#1}{o}{%
    \par\Needspace{4\baselineskip}\addvspace{\medskipamount}\noindent
    \textbf{#2}\IfValueT{##1}{ (##1)}\textbf{.}\ \ignorespaces}%
  {\par\addvspace{\medskipamount}}}
\newrunin{esempio}{Esempio}
\newrunin{esercizio}{Esercizio}
\newrunin{osservazione}{Osservazione}

% V e F nelle tavole di verità
\newcommand{\V}{\textsf{V}}
\newcommand{\F}{\textsf{F}}

% --- Codice C ---
% Uso: \begin{codice} ... \end{codice}      codice C numerato
%      \begin{comando} ... \end{comando}    comandi da terminale
%      \begin{risultato} ... \end{risultato} output di un programma
%      \lstinline|printf("ciao");|           codice dentro al testo
\lstdefinestyle{c}{
  language=C,
  basicstyle=\ttfamily\small,
  keywordstyle=\color{blue}\bfseries,
  commentstyle=\color{gray}\itshape,
  stringstyle=\color{red!70!black},
  numbers=left,
  numberstyle=\tiny\color{gray},
  numbersep=8pt,
  columns=flexible,
  keepspaces=true,
  breaklines=true,
  showstringspaces=false,
  tabsize=4,
  morekeywords={include, define},
  % lettere accentate nei commenti
  literate={à}{{\`a}}1 {è}{{\`e}}1 {é}{{\'e}}1 {ì}{{\`i}}1
           {ò}{{\`o}}1 {ù}{{\`u}}1 {È}{{\`E}}1 {À}{{\`A}}1
}
\lstset{style=c}
\newtcblisting{codice}{listing only, nota, left=6mm,
  listing options={style=c}}
\newtcblisting{comando}{listing only, nota,
  listing options={basicstyle=\ttfamily\small, numbers=none, language={},
    columns=flexible, keepspaces=true, breaklines=true}}
\newtcblisting{risultato}{listing only, enhanced jigsaw, breakable,
  colback=black!5, colframe=black, boxrule=0.5pt, sharp corners,
  listing options={basicstyle=\ttfamily\small, numbers=none, language={},
    columns=flexible, keepspaces=true, breaklines=true}}

% --- Diagrammi di flusso (simboli standard) ---
\tikzset{
  terminale/.style = {ellipse, draw, minimum width=2.5cm, minimum height=0.9cm},
  io/.style        = {trapezium, trapezium left angle=70, trapezium right angle=110,
                      draw, minimum width=2.5cm, minimum height=0.9cm},
  processo/.style  = {rectangle, draw, minimum width=2.5cm, minimum height=0.9cm},
  decisione/.style = {diamond, draw, aspect=2, minimum width=2.5cm},
  freccia/.style   = {thick, -{Stealth}}
}

\title{Continuità delle funzioni elementari}
\author{Marco Cerbella}
\date{Lezione 11 -- 29 ottobre 2026}

\begin{document}
\maketitle

\section{Continuità delle funzioni elementari}

Le funzioni elementari:
\begin{itemize}
    \item funzioni potenza,
    \item funzioni esponenziali,
    \item funzioni logaritmiche,
    \item funzioni trigonometriche e loro inverse,
\end{itemize}
sono continue nel loro dominio.

\section{Verifiche di continuità con la definizione}

\begin{esempio}[$f(x) = ax + b$ è continua in ogni $x_0 \in \mathbb{R}$, con $a, b \in \mathbb{R}$]
Fissato $\varepsilon > 0$ consideriamo la disuguaglianza $|f(x) - f(x_0)| < \varepsilon$:
\[
|f(x) - f(x_0)| = |ax + b - (ax_0 + b)| = |ax - ax_0| = |a|\cdot|x - x_0| .
\]
Vogliamo trovare $\delta > 0$ per cui se $|x - x_0| < \delta$ allora $|a|\cdot|x - x_0| < \varepsilon$.
\begin{itemize}
    \item Se $a = 0$: $|a|\cdot|x-x_0| = 0 < \varepsilon$ per ogni $x$.
    \item Se $a \neq 0$: $|a|\cdot|x - x_0| < \varepsilon \iff |x - x_0| < \dfrac{\varepsilon}{|a|}$; si sceglie $\delta = \dfrac{\varepsilon}{|a|}$ e la definizione è verificata.
\end{itemize}
\end{esempio}

\begin{esempio}[$f(x) = x^2$ è continua in $x_0 = 2$]
Fissato $\varepsilon > 0$:
\[
|f(x) - f(x_0)| = |x^2 - 4| = |x - 2|\cdot|x + 2| ,
\]
usando $a^2 - b^2 = (a - b)(a + b)$ per ogni $a, b \in \mathbb{R}$. Possiamo considerare $x \in I_1(2)$, cioè $2 - 1 < x < 2 + 1$; allora $1 < x < 3$ e quindi $|x + 2| < 5$. Dunque
\[
|x - 2|\,|x + 2| < |x - 2| \cdot 5 < \varepsilon \iff |x - 2| < \frac{\varepsilon}{5} .
\]
Consideriamo $\delta = \min\{1, \frac{\varepsilon}{5}\}$ e abbiamo che la definizione di continuità è soddisfatta, quindi $f(x) = x^2$ è continua in $x_0 = 2$. In realtà $f(x) = x^2$ è continua per ogni $x \in \mathbb{R}$.
\end{esempio}

\begin{esempio}[$\lim_{x\to 1}(2x^2 + 3) = 5$]
Qui $f(x) = 2x^2 + 3$, $x_0 = 1$, $\ell = 2\cdot 1 + 3 = 5$. Verifichiamo con la definizione. Fissato $\varepsilon > 0$ consideriamo $|f(x) - \ell| < \varepsilon$:
\[
|2x^2 + 3 - 5| = |2x^2 - 2| = 2|x^2 - 1| = 2|x - 1|\cdot|x + 1| .
\]
Possiamo supporre $x \in I_1(1)$, cioè $1 - 1 < x < 1 + 1$, $0 < x < 2$; allora $|x + 1| < 2 + 1 = 3$, e
\[
2|x - 1|\,|x + 1| < 6\,|x - 1| < \varepsilon \quad\text{se}\quad |x - 1| < \delta = \frac{\varepsilon}{6}
\]
(più precisamente $\delta = \min\{1, \varepsilon/6\}$).
\end{esempio}

\section{Algebra delle funzioni continue}

\begin{teorema}[Algebra delle funzioni continue]
Se $f$ e $g$ sono continue in $x_0$, allora $f + g$, $f \cdot g$ e $f/g$ (se $g(x_0) \neq 0$) sono continue in $x_0$.
\end{teorema}

\begin{teorema}[Composizione di funzioni continue]
Se $g \circ f$ esiste, e $f$ è continua in $x_0$ e $g$ è continua in $f(x_0)$, allora la funzione composta $g \circ f$ è continua in $x_0$.
\end{teorema}

\begin{osservazione}[Conseguenza]
Una funzione espressa come numero finito di operazioni elementari è continua. Per calcolare il limite di una tale funzione è sufficiente calcolare la funzione in quel punto.
\end{osservazione}

\begin{esempio}[Determinare per quali $\alpha \in \mathbb{R}$ la funzione è continua]
\[
f(x) = \begin{cases} \alpha \cdot \sin\!\left(x + \dfrac{\pi}{2}\right) & x > 0 \\[2mm] 2x^2 + 3 & x \le 0 .\end{cases}
\]
Per $x < 0$ e per $x > 0$ la funzione è continua perché somma e prodotto di funzioni elementari continue. Verifichiamo la continuità in $x = 0$:
\[
\lim_{x\to 0^+} \alpha \sin\!\left(x + \frac{\pi}{2}\right) = \alpha \sin\frac{\pi}{2} = \alpha, \qquad
\lim_{x\to 0^-} (2x^2 + 3) = 3, \qquad f(0) = 3 .
\]
Perché $\lim_{x\to 0^+} f(x) = \lim_{x\to 0^-} f(x) = f(0)$ bisogna che $\alpha = 3$. Dunque la funzione $f$ è continua in $\mathbb{R}$ per $\alpha = 3$.
\end{esempio}

\end{document}
